% problem-000044 4 铁钼催化剂EDTA分析
\section*{4. 铁钼催化剂EDTA分析}
\textit{下列为分问。}

\subsection*{4-1}
$\mathrm { Z n ^ { 2 + } }$ 标准溶液的配制：称取0.3289 g⾦属 $Z \mathrm { n }$ ，溶解后转移⾄250 mL的容量瓶，定容。EDTA溶液的标定：⽤移液管取25.00 mL 上述 $\mathrm { Z n ^ { 2 + } }$ 标准溶液，以⼆甲酚橙作指⽰剂，⽤EDTA标准溶液滴⾄终点，消耗20.24 mL。计算 $\mathrm { Z n ^ { 2 + } }$ 标准溶液和 EDTA 标准溶液的浓度。

\subsection*{4-2}
实验误差⽤标准差表⽰。对�次测量结果作加减和乘除运算，运算结果(�)的误差为 ${ \boldsymbol { \Delta Q } }$ ，其计算公式分别为 $\Delta Q _ { \pm } = \sqrt { \sum _ { i = 1 } ^ { n } ( \Delta q _ { i } ) ^ { 2 } } \mathcal { \vec { H } } \mathbb { I } \Delta Q _ { \times \div } = | Q | \sqrt { \sum _ { i = 1 } ^ { n } ( \Delta q _ { i } / q _ { i } ) ^ { 2 } }$ ，其中 $\boldsymbol { \cdot } \Delta \boldsymbol { q } _ { i }$ 为第�次测量值 $\cdot q _ { i }$ 的误差。已知分析天平称量的标准差为0.1 mg，移液管、容量瓶和滴定管的标准差均为0.01 mL，忽略其他误差因素，计算4-1中EDTA标准溶液浓度的最⼩标准差。

\subsection*{4-3}
请计算该催化剂中Fe和Mo的原⼦个数之⽐。

\subsection*{4-4}
欲在室温下配制1.0 L pH 5.00的醋酸-醋酸钠缓冲溶液，现有11.8 g醋酸钠，请计算需要加⼊醋酸 $( \rho =$ $1 . 0 5 \mathrm { g } \mathrm { m L } ^ { - 1 } )$ 的体积（醋酸在⽔溶液中的 $K _ { a } _ { \mathrm { ~ } } = 1 . 8 \times 1 0 ^ { - 5 } )$ 0
